\chapter{Exercices progressifs : faire sans regarder les corrigés}
\textbf{Consigne commune.} Donner des valeurs exactes. Les angles sans unité précisée sont en radians. Lorsque \(k\in\Z\), préciser le paramètre entier. Pour les équations sur un intervalle, fournir l'ensemble complet des solutions.

\exercice{1}{Seconde — conversions et cercle}
Convertir \(225^\circ\) en radians et \(-7\pi/6\) en degrés. Pour un cercle de rayon \(6\) cm et un secteur d'angle \(5\pi/6\), déterminer la longueur de l'arc puis l'aire du secteur.

\exercice{2}{Seconde — angles remarquables}
Calculer exactement \(\sin\theta\) et \(\cos\theta\) pour les trois angles : \(\theta=11\pi/6\), \(\theta=-7\pi/4\), \(\theta=17\pi/3\).

\exercice{3}{Seconde — signe et Pythagore}
On sait que \(\cos x=4/5\) et que \(x\in]3\pi/2,2\pi[\). Calculer \(\sin x\) et \(\tan x\).

\exercice{4}{Seconde — triangle rectangle}
Dans un triangle rectangle, un angle aigu vaut \(\pi/6\) et l'hypoténuse mesure \(14\) cm. Donner les longueurs des côtés adjacent et opposé.

\exercice{5}{Première — identités et domaines}
Pour tout réel \(x\), justifier \(\sin(\pi-x)=\sin x\) et \(\cos(x+\pi)=-\cos x\). Simplifier \(\displaystyle\frac{\sin(2x)}{2\cos x}\), en précisant le domaine sur lequel l'expression initiale est définie.

\exercice{6}{Première — fonction sinusoïdale}
Soit \(f(x)=3\sin(2x+\pi/6)-2\). Déterminer son amplitude, sa période fondamentale, son décalage horizontal, sa valeur centrale et son image. Calculer \(f'(x)\).

\exercice{7}{Première — formules d'addition}
Calculer exactement \(\sin(\pi/12)\) et \(\cos(\pi/12)\) à l'aide de \(\pi/3-\pi/4\).

\exercice{8}{Première — linéarisation}
Démontrer \(\cos^2x=(1+\cos2x)/2\). En déduire une primitive de \(\cos^2x\).

\exercice{9}{Seconde / Première — équation}
Résoudre \(\sin x=-1/2\) dans \([0,2\pi[\).

\exercice{10}{Première — équation avec argument double}
Résoudre \(\cos(2x)=1/2\) dans \([0,2\pi[\).

\exercice{11}{Première — factorisation}
Résoudre \(\sin(2x)=\cos x\) dans \([0,2\pi[\). Expliquer pourquoi il serait dangereux de diviser immédiatement par \(\cos x\).

\exercice{12}{Première — inéquation}
Résoudre \(\sin x\geq\sqrt2/2\) dans \([0,2\pi]\), puis exprimer l'ensemble des solutions dans \(\R\).

\exercice{13}{Première / Terminale — combinaison trigonométrique}
Résoudre dans \(\R\) l'équation \(\sin x+\cos x=1\). Donner ensuite les solutions dans \([0,2\pi[\).

\exercice{14}{Première — triangle quelconque}
Dans un triangle \(ABC\), \(b=5\), \(c=7\), \(A=\pi/3\). Déterminer exactement le côté \(a\) et l'aire du triangle.

\exercice{15}{Première — produit scalaire}
Pour \(\vec u=(1,2)\) et \(\vec v=(2,-1)\), calculer leur produit scalaire, leurs normes et l'angle non orienté entre ces vecteurs.

\exercice{16}{Terminale — phénomène périodique}
La hauteur d'un point en mètres est \(h(t)=2+1{,}5\sin(\pi t/6)\), pour \(t\geq0\) en secondes. Trouver la période, la fréquence, le minimum, le maximum et le premier instant de hauteur maximale.

\exercice{17}{Terminale — variations}
Étudier les extrema de \(f(x)=\sin x+\cos x\) sur \([0,2\pi]\), à l'aide de la dérivée. Vérifier les extrema à l'aide d'une identité trigonométrique.

\exercice{18}{Terminale / Approfondissement — réciproques}
Calculer exactement \(\arcsin(\sin(5\pi/6))\), \(\arccos(\cos(5\pi/3))\) et \(\arctan(\tan(3\pi/4))\).

\exercice{19}{Première — tangente}
Résoudre \(\tan x=\sqrt3\) dans \(\R\), puis dans \([0,2\pi[\). Préciser l'ensemble de définition de la tangente.

\exercice{20}{Première — identité et équation}
Résoudre \(\cos(2x)=\cos x\) dans \([0,2\pi[\) de deux manières : avec la formule de l'angle double, puis avec l'équivalence \(\cos a=\cos b\).
